\documentclass[10pt]{article}
\usepackage{icmlko}

\icmlmeta{IP 파운데이션 월드모델}{312}{확실한 것은 도메인 하나다}{잡음이 작은 자로 빌림 표를 다시 쟀다 --- 도서는 굳어지고 펀딩은 무너진다}

\begin{document}

\twocolumn[
\icmlheader

\begin{abstract}
노트 311이 ``순위를 정하는 자와 현상을 재는 자는 같을 이유가 없다''로
끝났다. F6 의 짝SE 가 F18 의 3분의 2다. 노트 303의 빌림 표 전체를 F6 으로
다시 쟀다. \textbf{도서가 굳어지고 펀딩이 무너진다.} 도서 $+$0.1143 ·
짝SE 0.0302 · \textbf{$t{=}{-}3.79$} --- 도메인 문턱 2.77 을 넉넉히 넘는다.
펀딩은 F18 에서 0.1089(43\%)였던 것이 F6 에서 \textbf{0.0474}(29\% ·
$t{=}{-}1.53$)로 준다. 나머지 여섯은 $-$0.0242 $\sim$ $+$0.0017 로 0 이거나
살짝 음수다(막으면 조금 나아진다). 정식화 넷을 나란히 놓으면 더 뚜렷하다 ---
\textbf{도서는 F6 3.79 · F9 3.45 · F21 3.30 으로 셋에서 문턱을 넘고}
펀딩은 F6 1.53 · F18 2.00 으로 \textbf{어디서도 못 넘는다.} 그래서 이
실험실의 파운데이션 효과를 정직하게 적으면 \textbf{도메인 하나 · 유보
163행 · 판의 6\%}다. 노트 303이 ``둘 · 243행 · 9\%''로 적은 것을 줄인다.
\textbf{줄었지만 굳었다} --- 그전까지 이 현상은 어느 자로 재도 문턱 언저리에
있었고 이제 하나가 확실하다.
\end{abstract}
]

\section{자를 바꿔 다시 쟀다}

노트 311이 정식화 여섯으로 도서의 빌림을 재고 \textbf{선형의 짝SE 가
나무의 3분의 2}인 것을 봤다(0.030$\sim$0.034 대 0.042$\sim$0.048). 같은
현상을 더 또렷하게 보는 자가 있다는 뜻이다. 노트 303의 표 전체를 F6 으로
다시 쟀다.

\begin{figure*}[t]
\centering
\includegraphics[width=\textwidth]{figs/one.pdf}
\caption{왼쪽 --- F6 으로 잰 도메인별 빌림($\pm$2$\times$짝SE, 부트스트랩
400, 괄호는 유보 행). 오른쪽 --- 도서와 펀딩을 정식화 넷에서.}
\end{figure*}

\begin{center}\footnotesize
\begin{tabular}{lrrrrr}
\hline
도메인 & 유보 & 빌림 & 짝SE & $t$ & 몫 \\
\hline
\textbf{도서} & 163 & \textbf{$+$0.1143} & 0.0302 & \textbf{$-$3.79} & 37\% \\
펀딩 & 80 & $+$0.0474 & 0.0311 & $-$1.53 & 29\% \\
애니 & 606 & $+$0.0017 & 0.0062 & $-$0.26 & 0\% \\
웹툰 & 711 & $-$0.0019 & 0.0040 & $+$0.47 & 0\% \\
게임 & 180 & $-$0.0027 & 0.0057 & $+$0.47 & 0\% \\
세계애니 & 300 & $-$0.0051 & 0.0150 & $+$0.34 & $-$1\% \\
모바일 & 441 & $-$0.0069 & 0.0036 & $+$1.91 & $-$1\% \\
아이돌 & 25 & $-$0.0242 & 0.0338 & $+$0.72 & $-$9\% \\
\hline
\end{tabular}
\end{center}

\section{도서는 굳고 펀딩은 무너진다}

\begin{center}\small
\begin{tabular}{lrr}
\hline
정식화 & 도서 $|t|$ & 펀딩 $|t|$ \\
\hline
\textbf{F6} (선형) & \textbf{3.79} & 1.53 \\
\textbf{F9} (선형) & \textbf{3.45} & --- \\
\textbf{F21} (선형) & \textbf{3.30} & --- \\
F18 (나무) & 2.64 & 2.00 \\
\hline
\end{tabular}
\end{center}

\textbf{도서는 셋에서 문턱(2.77)을 넘는다.} 펀딩은 나무에서 2.00, 선형에서
1.53 --- \textbf{어디서도 못 넘는다.} 노트 303이 F18 로 재서 ``펀딩 44\%''를
적었는데 그 수는 \textbf{자를 바꾸면 29\%로 준다.} 유보 80행이라 짝SE 가
0.031 이고 그 크기에서는 못 가른다.

\paragraph{나머지 여섯은 확실히 0 이다} 이번에는 짝SE 가 0.0036$\sim$0.0150
으로 작아서 ``못 잰다''가 아니라 \textbf{정말 0}이라고 말할 수 있다.
모바일 $-$0.0069 에 짝SE 0.0036 이면 $\pm$0.007 안에 있다는 뜻이다.
\textbf{큰 도메인이 안 빌리는 것이 표본 부족 때문이 아니다.}

\section{그래서 크기를 정직하게 적는다}

\begin{center}\small
\begin{tabular}{lrr}
\hline
& 노트 303 & \textbf{이 노트} \\
\hline
빌리는 도메인 & 둘 & \textbf{하나} \\
그 유보 행 & 243 & \textbf{163} \\
판에서의 몫 & 9\% & \textbf{6\%} \\
$|t|$ & 2.87 · 1.97 & \textbf{3.79} \\
\hline
\end{tabular}
\end{center}

\textbf{줄었지만 굳었다.} 노트 303 · 309 · 311을 거치면서 이 현상은 계속
문턱 언저리에 있었다($-$2.87 $\to$ $-$2.54 $\to$ $-$2.64). 자를 바꾸니
\textbf{$-$3.79} 다.

\paragraph{이것이 이 실험실의 파운데이션 효과다} 열한 도메인 · 유보
2,669행 · 축 스물넷을 쌓아 온 판에서, \textbf{도메인 가로지르기가 확실하게
일하는 자리는 도서 163행이다.} 그 163행에서 제 점수의 \textbf{37\%}가
남에게서 온다. 나머지 2,506행은 자기 자료로 돈다.

\paragraph{작다고 무의미하지 않다} 도서는 학습이 80행뿐이다 --- 이 판에서
팝업(16) 다음으로 얇다. \textbf{얇은 도메인이 두꺼운 도메인의 구조를 받아
$\rho$ 0.31 을 낸다는 것}이 파운데이션이 하겠다는 일 그대로다. 다만
\textbf{지금은 그런 도메인이 하나뿐이다.}

\paragraph{남은 것}
\begin{center}\small
\begin{tabular}{ll}
\hline
갈래 & 왜 \\
\hline
얇은 도메인을 더 & 빌릴 자리를 늘리는 유일한 길 \\
펀딩 유보 80행 & 문턱 0.06 이 너무 높다 \\
도서에 축을 더 & \texttt{publisher} 는 떨어졌다 \\
시장층 위약 & 노트 304의 셋째 층 \\
\hline
\end{tabular}
\end{center}

\paragraph{규약} \textbf{현상이 확실해지면 크기도 같이 줄어드는 것이
보통이다.} 노트 303이 ``둘 · 9\%''로 적었을 때 그 둘 다 문턱 아래였다.
자를 바꾸니 하나는 넘고 하나는 더 내려갔다 --- \textbf{합쳐서 크게 보이던
것이 갈리면 작아진다}(노트 298에서 하네스 셋을 합쳤을 때와 같다).
\textbf{그리고 0 을 0 이라고 말하려면 짝SE 가 작아야 한다} --- 나머지 여섯이
0 인 것을 F18 로는 ``못 잰다''로밖에 못 적었는데 F6 에서는 $\pm$0.007 이다.

\end{document}
